% ----------------------------------------------------------------------
%  3.5.7. Сверхпроводимость
% ----------------------------------------------------------------------
\subsection{Сверхпроводимость}\label{ssec:supercond}

В 1911 г. Камерлинг-Оннес, измеряя сопротивление ртути при низких
температурах, обнаружил, что при 4,15~К оно \emph{скачком падает до
нуля} (рис.~\ref{fig:sc-basic},\,\textit{а}). Температуру перехода
называют \term{критической} $T_\text{к}$.

\paragraph{Где встречается}
\begin{itemize}
\item Более 40 металлических элементов в тех или иных условиях становятся
сверхпроводниками. Среди чистых элементов самая высокая $T_\text{к}$ у
ниобия~— 9,2~К; очень низкие~— у вольфрама (0,015~К)\ispr{В оригинале
(с.~513) для вольфрама указано $T_\text{к}=0{,}001$~К; измеренное значение~—
0,015~К.}.
\item Соединения со структурой A15: Nb$_3$Sn~— 18~К, Nb$_3$Ge~— 23,3~К.
\item В 1986 г. Беднорц и Мюллер открыли \term{высокотемпературную
сверхпроводимость} (ВТСП) в системе BaO--La$_2$O$_3$--CuO ($T_\text{к}=30$--35~К).
Теперь известны оксидные сверхпроводники с $T_\text{к}>100$~К: перейдён
технически важный порог~— температура кипения жидкого азота (77~К),
и жидкий гелий (4,2~К) для них не нужен.
\end{itemize}

\paragraph{Критическое поле и критический ток}
Сверхпроводимость разрушается не только нагревом, но и магнитным
полем выше \term{критического} $H_\text{к}$, и током выше
\term{критической плотности} $j_\text{к}$. Критическое поле растёт при
охлаждении (рис.~\ref{fig:sc-basic},\,\textit{б}):
\begin{keyeq}\label{eq:Hc}
H_\text{к}(T)=H_\text{к}(0)\left[1-\left(\frac T{T_\text{к}}\right)^2\right].
\end{keyeq}
Простота этого закона позволяет находить большие поля при $T\to0$,
экстраполируя измерения в слабых полях. Рекордные значения: для
низкотемпературных сверхпроводников $B_\text{к}=\mu_0H_\text{к}\approx60$~Тл
(фаза Шевреля PbMo$_6$S$_8$), у керамических ВТСП ещё больше;
$j_\text{к}$ до $10^{11}$~А/м$^2$ при $T\to0$ (соединения A15).

\paragraph{Эффект Мейсснера}
Мейсснер и Оксенфельд обнаружили, что при охлаждении в магнитном поле
сверхпроводник при переходе \emph{выталкивает} поле из себя: внутри
$B=0$ (рис.~\ref{fig:sc-basic},\,\textit{в}). Сверхпроводник~—
\term{идеальный диамагнетик}. Это такое же фундаментальное его
свойство, как нулевое сопротивление. Поле всё же проникает в
поверхностный слой глубиной $\lambda\sim10^{-8}$~м (\term{глубина
проникновения}); в нём текут незатухающие токи, магнитный момент
которых компенсирует внешнее поле. При переходе скачком меняется
электронная теплоёмкость, а решёточная и сама кристаллическая
структура не меняются.

\begin{figure}[ht]
\centering
\begin{subfigure}[b]{0.31\textwidth}\centering
\begin{tikzpicture}
\begin{axis}[width=1.05\textwidth, height=4.4cm, xmin=3.9, xmax=4.5, ymin=0, ymax=0.16,
  xlabel={$T$, К}, ylabel={$R$, Ом}, axis lines=left, xtick={4.0,4.2,4.4},
  xticklabel style={/pgf/number format/use comma}, ytick=\empty]
  \addplot[fmBlue, thick] coordinates {(3.9,0) (4.15,0) (4.155,0.11) (4.5,0.125)};
  \node[font=\scriptsize, anchor=west] at (axis cs:4.17,0.05) {$T_\text{к}$};
\end{axis}
\end{tikzpicture}
\caption{}
\end{subfigure}\hfill
\begin{subfigure}[b]{0.31\textwidth}\centering
\begin{tikzpicture}
\begin{axis}[width=1.05\textwidth, height=4.4cm, xmin=0, xmax=1.1, ymin=0, ymax=1.15,
  xlabel={$T/T_\text{к}$}, ylabel={$H_\text{к}/H_\text{к}(0)$}, axis lines=left, xtick={0,1}, ytick={0,1}]
  \addplot[fill=fmBlue!15, draw=fmBlue, thick, domain=0:1, samples=40] {1-x^2} \closedcycle;
  \node[font=\scriptsize] at (axis cs:0.38,0.4) {сверхпров.};
  \node[font=\scriptsize] at (axis cs:0.82,0.88) {норм.};
\end{axis}
\end{tikzpicture}
\caption{}
\end{subfigure}\hfill
\begin{subfigure}[b]{0.31\textwidth}\centering
\begin{tikzpicture}[scale=0.8]
  \fill[fmBlue!20] (0,0) circle (0.8);
  \draw[fmBlue, thick] (0,0) circle (0.8);
  \foreach \y in {-1.3,-0.9,-0.45,0.45,0.9,1.3}{
    \pgfmathsetmacro{\yy}{\y*(1+0.6*exp(-3*abs(\y)+1.2))}
    \draw[fmRed, thick, smooth] (-2,\y) .. controls (-1,\y) and (-0.8,\yy) .. (0,\yy) .. controls (0.8,\yy) and (1,\y) .. (2,\y);}
  \node[font=\scriptsize] at (0,0) {$B=0$};
\end{tikzpicture}
\caption{}
\end{subfigure}
\caption{\textit{а}~— исчезновение сопротивления ртути (опыт
Камерлинг-Оннеса, схема); \textit{б}~— граница сверхпроводящего
состояния на диаграмме $H$--$T$; \textit{в}~— эффект Мейсснера:
сверхпроводник выталкивает силовые линии магнитного поля}
\label{fig:sc-basic}
\end{figure}

\paragraph{Природа сверхпроводимости: куперовские пары}
Первая феноменологическая теория (Ф. и Г.~Лондоны) представляла
сверхпроводник как смесь нормальных и «сверхпроводящих» электронов, но
не объясняла, что это за электроны. Микроскопическую теорию построили
Бардин, Купер, Шриффер и Н.\,Н.~Боголюбов (\term{теория БКШ}).

Электроны~— фермионы, и для системы фермионов сверхпроводимость (как и
сверхтекучесть, открытая у жидкого гелия ниже 2,17~К) невозможна: такие
явления свойственны бозонам~— частицам с целым спином. Значит, нужно
найти механизм, связывающий электроны в \emph{пары} с нулевым спином.
Как могут притягиваться одноимённо заряженные частицы?
\begin{shagi}
\item Электрон, пролетая между ионами, притягивает их, и они начинают
смещаться к его траектории. Ионы инертны: скорость электрона около
$10^6$~м/с, период колебаний ионов около $10^{-13}$~с, и пока ионы
сместятся, электрон уходит на $\sim10^{-7}$~м\ispr{В оригинале (с.~516)
скорость электрона указана как $\sim10^6$~см/с; при такой скорости
электрон сместился бы лишь на $\sim10^{-9}$~м. Скорость фермиевских
электронов~— $\sim10^6$~м/с.}.
\item За электроном остаётся след~— область избыточного положительного
заряда, к которой притягивается второй электрон.
\item Если это притяжение сильнее кулоновского отталкивания, возникает
связанное состояние~— \term{куперовская пара}.
\end{shagi}

На квантовом языке электроны обмениваются фононами: первый испускает
фонон ($\vb p_1=\vb p_1'+\vb q$), второй его поглощает
($\vb p_2+\vb q=\vb p_2'$); суммарный импульс сохраняется. Энергия системы
понижается сильнее всего, когда разрешено наибольшее число таких
процессов рассеяния~— для пар с \emph{противоположными импульсами и
спинами}. Спариваются не все электроны, а лишь в тонком слое у
поверхности Ферми.

\begin{glavnoe}
Куперовская пара имеет нулевой спин и подчиняется статистике
Бозе--Эйнштейна. Все пары «конденсируются» в одно квантовое состояние с
общей волновой функцией и фазой. Ток в сверхпроводнике~— движение этой
единой волны, которую решётка не рассеивает: фононы сами участвуют в
её образовании. Чтобы рассеять пару, её нужно разрушить, а для этого
нужна конечная энергия.
\end{glavnoe}

\paragraph{Энергетическая щель и длина когерентности}
\begin{itemize}
\item Разрушение пары требует энергии не меньше $\Delta$~— в спектре
появляется \term{энергетическая щель}. По теории БКШ
$\Delta(0)=3{,}5\,kT_\text{к}$, а вблизи $T_\text{к}$ щель исчезает:
$\Delta(T)=1{,}74\,\Delta(0)\sqrt{1-T/T_\text{к}}$.
\item Размер пары~— \term{длина когерентности} $\xi\sim10^{-7}$--$10^{-6}$~м.
В объёме одной пары находятся центры миллионов других пар: пары не
независимы, спаривание~— коллективный эффект.
\end{itemize}
Теория БКШ не описывает высокотемпературную сверхпроводимость:
установлено, что её механизм не фононный. Предлагались экситонный
механизм, модель резонирующих валентных связей и др., но окончательной
теории ВТСП пока нет.

\paragraph{Сверхпроводники I и II рода}
\begin{itemize}
\item \emph{I рода}~— чистые металлы (кроме Nb, V, Tc): малые
критические поля, полное выталкивание поля вплоть до $H_\text{к}$.
\item \emph{II рода}~— в основном сплавы и соединения. У них два
критических поля (рис.~\ref{fig:sc-types}). Ниже $H_\text{к1}$~— полный
эффект Мейсснера. При $H_\text{к1}<H<H_\text{к2}$ поле частично проникает
в образец, а сопротивление остаётся нулевым~— \term{смешанное
состояние}: в сверхпроводящей матрице появляются нормальные нити
диаметром около $2\xi$ вдоль поля~— \term{вихри}, каждый несёт один
\emph{квант магнитного потока} $\Phi_0=h/2e=\pi\hbar/e$. В монокристалле
вихри образуют правильную решётку. С ростом поля их становится больше,
и при $H_\text{к2}$ они смыкаются~— сверхпроводимость
исчезает\ispr{В оригинале (с.~519) $B_\text{к2}=\Phi_0/\pi\xi^2=h/e\xi^2$;
второе равенство не согласуется с первым ($\Phi_0=h/2e$), а теория
Гинзбурга--Ландау даёт $B_\text{к2}=\Phi_0/(2\pi\xi^2)$.}:
\end{itemize}
\begin{equation}\label{eq:Hc2}
B_\text{к2}=\frac{\Phi_0}{2\pi\xi^2}.
\end{equation}
Чем меньше длина
когерентности, тем выше верхнее критическое поле.

Теория Гинзбурга--Ландау--Абрикосова--Горькова (ГЛАГ) показала, что тип
сверхпроводника определяется параметром $\varkappa=\lambda/\xi$: при
$\varkappa<1/\sqrt2$~— I род, при $\varkappa>1/\sqrt2$~— II род. При
малой длине когерентности и большой глубине проникновения энергия
границы между нормальной и сверхпроводящей фазами отрицательна, и
выгодно смешанное состояние.

\begin{figure}[ht]
\centering
\begin{subfigure}[b]{0.48\textwidth}\centering
\begin{tikzpicture}
\begin{axis}[width=\textwidth, height=4.6cm, xmin=0, xmax=3.3, ymin=0, ymax=1.25,
  xlabel={$H$}, ylabel={$-M$}, xtick={1,1.6,3}, xticklabels={$H_\text{к1}$,$H_\text{к}$,$H_\text{к2}$},
  ytick=\empty, axis lines=left]
  \addplot[fmGray, thick, dashed] coordinates {(0,0) (1.6,1.6*0.7) (1.6,0)};
  \addplot[fmBlue, thick, smooth] coordinates {(0,0) (1,0.7)};
  \addplot[fmBlue, thick, smooth] coordinates {(1,0.7) (1.15,0.5) (1.5,0.33) (2,0.19) (2.5,0.08) (3,0)};
  \node[font=\scriptsize, fmGray, anchor=south] at (axis cs:1.4,1.15) {I род};
  \node[font=\scriptsize, fmBlue, anchor=west] at (axis cs:1.9,0.35) {II род};
\end{axis}
\end{tikzpicture}
\caption{}
\end{subfigure}\hfill
\begin{subfigure}[b]{0.48\textwidth}\centering
\begin{tikzpicture}[scale=0.62]
  \fill[fmBlue!15] (0,0) rectangle (5.2,4.2);
  \draw[fmBlue] (0,0) rectangle (5.2,4.2);
  \foreach \i in {0,...,4}{\foreach \j in {0,...,3}{
    \pgfmathsetmacro{\xx}{0.6+\i*1.0+mod(\j,2)*0.5}
    \pgfmathsetmacro{\yy}{0.55+\j*0.95}
    \fill[fmRed!25] (\xx,\yy) circle (0.27);
    \draw[fmRed, thick] (\xx,\yy) circle (0.1);
    \fill[fmRed] (\xx,\yy) circle (0.04);}}
  \node[font=\scriptsize, fill=white, inner sep=1pt] at (2.6,-0.45) {вихри: поле $\perp$ рисунку};
\end{tikzpicture}
\caption{}
\end{subfigure}
\caption{Сверхпроводники II рода: \textit{а}~— кривые намагничивания I
рода (пунктир) и II рода (сплошная линия, схема); \textit{б}~— решётка
вихрей в смешанном состоянии: в каждой нормальной сердцевине (диаметр
$\sim2\xi$) проходит один квант потока, вокруг текут
сверхпроводящие токи}
\label{fig:sc-types}
\end{figure}

\paragraph{Пиннинг и жёсткие сверхпроводники}
Когда ток течёт через сверхпроводник II рода в смешанном состоянии, на
вихри действует сила, перпендикулярная полю и току; движение вихрей
рассеивает энергию и может разрушить сверхпроводимость. Большие токи
возможны лишь при \emph{закреплённых} вихрях~— это \term{жёсткие}
сверхпроводники, а закрепление называют \term{пиннингом}. Центры
пиннинга~— скопления дислокаций, границы зёрен и особенно дисперсные
включения несверхпроводящих фаз. Здесь физика сверхпроводимости прямо
встречается с материаловедением: критический ток технических
сверхпроводников (Nb--Ti, Nb$_3$Sn) задаётся их микроструктурой. Кривые
намагничивания жёстких сверхпроводников сильно гистерезисны.

\paragraph{Эффект Джозефсона и ВТСП}
Джозефсон открыл туннелирование сверхпроводящих электронов через
тонкий (около 1~нм) слой изолятора между двумя сверхпроводниками. На
этом основаны \emph{сквиды}~— сверхпроводящие квантовые интерферометры,
измеряющие напряжения порядка $10^{-15}$~В и магнитные поля порядка
$10^{-15}$~Тл\footnote{В оригинале приведена чувствительность к полю
$\sim10^{-18}$~Тл; для серийных сквидов типичный уровень~— единицы
фемтотесла ($10^{-15}$~Тл).}.

Наиболее изучены ВТСП на основе иттрия YBa$_2$Cu$_3$O$_y$ ($y=6$--7, структура
перовскита с дефицитом кислорода, $T_\text{к}\approx92$~К) и висмута
(Bi--Me$^{2+}$--Cu--O). На монокристаллах, отдельных кристаллитах и тонких
плёнках при температуре жидкого азота получены плотности тока более
$10^9$~А/м$^2$ даже в сильных полях. Но в объёмных материалах (керамике,
проводах) критические токи гораздо меньше и сильно подавляются даже
слабыми полями. Главная причина~— ослабление сверхпроводимости на
контактах зёрен из-за очень малой длины когерентности ВТСП.

\begin{itogi}
\begin{itemize}
\item $\sigma=ne^2\tau/m_\text{эфф}$; в проводимости участвуют только
электроны у поверхности Ферми; эффективная масса задаётся кривизной
$E(k)$.
\item Заполненная зона тока не даёт; полупроводник~— заполненная
валентная зона и узкая щель; $n_i\propto T^{3/2}e^{-\Delta E/2kT}$;
доноры и акцепторы создают уровни в щели.
\item Правило Маттиссена: $\rho=\rho_\text{ост}+\rho_\text{ид}(T)$;
$\rho_\text{ид}\propto T$ при $T>\theta_D$ и $\propto T^5$ при $T\ll\theta_D$
(Блох--Грюнайзен).
\item Эффект Холла: $R_\text{Х}=1/ne$, знак~— знак носителей; положительный
знак у Be, Zn, Cd объясняется дырками.
\item Твёрдые растворы повышают сопротивление (парабола по составу),
упорядочение его снижает, двухфазные смеси дают аддитивную прямую.
\item Сверхпроводимость: $\rho=0$, эффект Мейсснера; куперовские пары,
щель $\Delta\approx3{,}5kT_\text{к}$; II род~— вихри и пиннинг.
\end{itemize}
\end{itogi}

\begin{proverka}
\item Почему модель Друде предсказывает $\rho\propto\sqrt T$, а опыт даёт
$\rho\propto T$?
\item Почему в электропроводности участвуют не все свободные электроны?
\item Чем отличается температурная зависимость проводимости металла и
полупроводника?
\item Как по отношению $R_{293}/R_{4{,}2}$ судить о чистоте металла?
\item Почему у бериллия постоянная Холла положительна?
\item Почему сопротивление сплава Al--Cu при старении при 20\,°C растёт, а при
225\,°C падает?
\item Почему для больших токов нужны сверхпроводники с пиннингом?
\end{proverka}
